\documentclass[11pt]{article}
\usepackage[margin=1in]{geometry}
\usepackage{amsmath,amssymb}
\usepackage{setspace}
\usepackage{url}
\onehalfspacing

\title{Project Idea: CG:SHOP 2027 Lawn Mowing}
\author{Hoang Ho, Kristina Liao, Oscar Martinez Wittinghan\\
COT 6405: Design and Analysis of Algorithms, Fall 2026}
\date{}

\begin{document}
\maketitle

\section{Problem statement}

This is the official CG:SHOP 2027 problem [1].

For a given polygonal region $P$, the Lawn Mowing Problem asks for a
shortest tour $T$ that gets within Euclidean distance $1/2$ of every
point in $P$. That is the same as a shortest tour for a unit-diameter
cutter $C$ that covers all of $P$. The problem is NP-hard [2].
CG:SHOP 2027 restricts instances to polyominoes. The cutter is a
smaller polyomino and moves in an axis-parallel fashion [1].

An instance gives a region to cover (a polyomino, possibly with holes),
a cutter shape, a cutter center, and an integer $k =$
\texttt{number\_of\_cutters}.

A solution is exactly $k$ tours, one per cutter.
Each tour lists the points the cutter's center passes through.
The tour is closed implicitly: the cutter goes from the last point back
to the first one.
Consecutive points must be axis-parallel.
A point may be visited again later, but not immediately again.
A tour of a single point is a cutter that stays where it is.
Cutters do not interfere with each other and may leave the region.
Together they must sweep the whole region.

The objective is $\mathrm{max\_tour\_length}$, the length of the
longest tour. We want to minimize it.

We will use the official checker from
\url{https://github.com/CG-SHOP/pyutils27} [3].
A solution counts only if that checker reports no errors.

\section{Goals}

We want work we can finish this semester, and still have enough for a
conference-style talk.

\begin{enumerate}
\item Implement more than one algorithm that produces feasible
      solutions (the official checker accepts them).
\item Run those algorithms on the published example instances, and on
      the challenge instances if they come out in time.
\item Compare them: $\mathrm{max\_tour\_length}$, runtime, and plots.
\item Give a talk with the problem, the algorithms, the plots, and
      what did not work.
\end{enumerate}

We are not trying to prove optimality, and we are not promising to win
the contest.

\section{Team and expected contributions}

\textbf{Hoang Ho, Kristina Liao, Oscar Martinez Wittinghan.}

We will share the coding, the experiments, and the writeup.
Who owns which algorithm is not fixed yet.
Everyone is expected to implement or test part of the solver, help
with the plots, and take part in the final talk.

\section{References}

\begin{enumerate}
\item S\'{a}ndor Fekete, Phillip Keldenich, and Stefan Schirra.
CG:SHOP 2027 Challenge.
\url{https://cgshop.ibr.cs.tu-bs.de/competition/cg-shop-2027/}.

\item Esther M. Arkin, S\'{a}ndor P. Fekete, and Joseph S. B. Mitchell.
Approximation algorithms for lawn mowing and milling.
{\em Computational Geometry}, 17(1--2):25--50, 2000.

\item CG:SHOP.
pyutils27.
\url{https://github.com/CG-SHOP/pyutils27}.
\end{enumerate}

\end{document}
